\documentclass[
  jou,
  floatsintext,
  longtable,
  a4paper,
  nolmodern,
  notxfonts,
  notimes,
  12pt,
  colorlinks=true,linkcolor=blue,citecolor=blue,urlcolor=blue]{apa7}

\usepackage{amsmath}
\usepackage{amssymb}



\usepackage[bidi=default]{babel}
\babelprovide[main,import]{spanish}


% get rid of language-specific shorthands (see #6817):
\let\LanguageShortHands\languageshorthands
\def\languageshorthands#1{}

\RequirePackage{longtable}
\RequirePackage{threeparttablex}

\makeatletter
\renewcommand{\paragraph}{\@startsection{paragraph}{4}{\parindent}%
	{0\baselineskip \@plus 0.2ex \@minus 0.2ex}%
	{-.5em}%
	{\normalfont\normalsize\bfseries\typesectitle}}

\renewcommand{\subparagraph}[1]{\@startsection{subparagraph}{5}{0.5em}%
	{0\baselineskip \@plus 0.2ex \@minus 0.2ex}%
	{-\z@\relax}%
	{\normalfont\normalsize\bfseries\itshape\hspace{\parindent}{#1}\textit{\addperi}}{\relax}}
\makeatother




\usepackage{longtable, booktabs, multirow, multicol, colortbl, hhline, caption, array, float, xpatch}
\usepackage{subcaption}


\renewcommand\thesubfigure{\Alph{subfigure}}
\setcounter{topnumber}{2}
\setcounter{bottomnumber}{2}
\setcounter{totalnumber}{4}
\renewcommand{\topfraction}{0.85}
\renewcommand{\bottomfraction}{0.85}
\renewcommand{\textfraction}{0.15}
\renewcommand{\floatpagefraction}{0.7}

\usepackage{tcolorbox}
\tcbuselibrary{listings,theorems, breakable, skins}
\usepackage{fontawesome5}

\definecolor{quarto-callout-color}{HTML}{909090}
\definecolor{quarto-callout-note-color}{HTML}{0758E5}
\definecolor{quarto-callout-important-color}{HTML}{CC1914}
\definecolor{quarto-callout-warning-color}{HTML}{EB9113}
\definecolor{quarto-callout-tip-color}{HTML}{00A047}
\definecolor{quarto-callout-caution-color}{HTML}{FC5300}
\definecolor{quarto-callout-color-frame}{HTML}{ACACAC}
\definecolor{quarto-callout-note-color-frame}{HTML}{4582EC}
\definecolor{quarto-callout-important-color-frame}{HTML}{D9534F}
\definecolor{quarto-callout-warning-color-frame}{HTML}{F0AD4E}
\definecolor{quarto-callout-tip-color-frame}{HTML}{02B875}
\definecolor{quarto-callout-caution-color-frame}{HTML}{FD7E14}

%\newlength\Oldarrayrulewidth
%\newlength\Oldtabcolsep


\usepackage{hyperref}




\providecommand{\tightlist}{%
  \setlength{\itemsep}{0pt}\setlength{\parskip}{0pt}}
\usepackage{longtable,booktabs,array}
\newcounter{none} % for unnumbered tables
\usepackage{calc} % for calculating minipage widths
% Correct order of tables after \paragraph or \subparagraph
\usepackage{etoolbox}
\makeatletter
\patchcmd\longtable{\par}{\if@noskipsec\mbox{}\fi\par}{}{}
\makeatother
% Allow footnotes in longtable head/foot
\IfFileExists{footnotehyper.sty}{\usepackage{footnotehyper}}{\usepackage{footnote}}
\makesavenoteenv{longtable}

\usepackage{graphicx}
\makeatletter
\newsavebox\pandoc@box
\newcommand*\pandocbounded[1]{% scales image to fit in text height/width
  \sbox\pandoc@box{#1}%
  \Gscale@div\@tempa{\textheight}{\dimexpr\ht\pandoc@box+\dp\pandoc@box\relax}%
  \Gscale@div\@tempb{\linewidth}{\wd\pandoc@box}%
  \ifdim\@tempb\p@<\@tempa\p@\let\@tempa\@tempb\fi% select the smaller of both
  \ifdim\@tempa\p@<\p@\scalebox{\@tempa}{\usebox\pandoc@box}%
  \else\usebox{\pandoc@box}%
  \fi%
}
% Set default figure placement to htbp
\def\fps@figure{htbp}
\makeatother


% definitions for citeproc citations
\NewDocumentCommand\citeproctext{}{}
\NewDocumentCommand\citeproc{mm}{%
  \begingroup\def\citeproctext{#2}\cite{#1}\endgroup}
\makeatletter
 % allow citations to break across lines
 \let\@cite@ofmt\@firstofone
 % avoid brackets around text for \cite:
 \def\@biblabel#1{}
 \def\@cite#1#2{{#1\if@tempswa , #2\fi}}
\makeatother
\newlength{\cslhangindent}
\setlength{\cslhangindent}{1.5em}
\newlength{\csllabelwidth}
\setlength{\csllabelwidth}{3em}
\newenvironment{CSLReferences}[2] % #1 hanging-indent, #2 entry-spacing
 {\begin{list}{}{%
  \setlength{\itemindent}{0pt}
  \setlength{\leftmargin}{0pt}
  \setlength{\parsep}{0pt}
  % turn on hanging indent if param 1 is 1
  \ifodd #1
   \setlength{\leftmargin}{\cslhangindent}
   \setlength{\itemindent}{-1\cslhangindent}
  \fi
  % set entry spacing
  \setlength{\itemsep}{#2\baselineskip}}}
 {\end{list}}
\usepackage{calc}
\newcommand{\CSLBlock}[1]{\hfill\break\parbox[t]{\linewidth}{\strut\ignorespaces#1\strut}}
\newcommand{\CSLLeftMargin}[1]{\parbox[t]{\csllabelwidth}{\strut#1\strut}}
\newcommand{\CSLRightInline}[1]{\parbox[t]{\linewidth - \csllabelwidth}{\strut#1\strut}}
\newcommand{\CSLIndent}[1]{\hspace{\cslhangindent}#1}





\usepackage{newtx}

\defaultfontfeatures{Scale=MatchLowercase}
\defaultfontfeatures[\rmfamily]{Ligatures=TeX,Scale=1}





\title{Institutos públicos de Ayacucho: sin luchas no hay victorias:
Cifras, normas y responsables: por qué la ley exige calidad sin asegurar
quién la paga, y una agenda estudiantil para exigirlo}


\shorttitle{INSTITUTOS PÚBLICOS DE AYACUCHO}


\usepackage{etoolbox}



\ccoppy{\textcopyright~2026}



\author{Edison Achalma}



\affiliation{
{Escuela Profesional de Economía, Universidad Nacional de San Cristóbal
de Huamanga}}




\leftheader{Achalma}

\date{2026-10-24}


\abstract{This essay examines why public technological institutes in
Ayacucho, Peru, struggle to meet the basic quality conditions required
for licensing, and turns a student manifesto into an argument supported
by official data and current law. It argues that the problem is not only
a lack of money but a design gap: Law 30512 requires every institute to
meet quality conditions, yet it guarantees resources for infrastructure
and equipment only to public teacher-training schools, while its reform
was enacted without additional funding. Official figures show the
consequences. Peru's public spending per tertiary student was 10.9
percent of GDP per capita in 2017, about half that of Chile or Colombia;
institutes and schools received 25.5 percent of the 2025 higher
education budget; and the Regional Government of Ayacucho, which runs
the institutes, assigned them 1.8 percent of its budget. Only 6 of the
16 public technological institutes of the education sector in the region
are licensed, while all its public teacher-training schools are. The
essay separates documented facts from interpretation, considers
alternative explanations such as execution capacity and territorial
dispersion, and reads the pattern through Mariátegui and Bourdieu. It
closes with a proposal, not yet a fact: a regional coordinating body of
institute students, with a first step, a tracking matrix of demands,
their legal basis, the authority responsible and a verifiable indicator
for each. }

\keywords{public higher education, technical and pedagogical
institutes, education spending, licensing, student
organization, Ayacucho}

\authornote{\par{\addORCIDlink{Edison Achalma}{0000-0001-6996-3364}} 

\par{   El autor no tiene conflictos de interés que revelar.  A los
estudiantes de institutos de Ayacucho que redactaron el pronunciamiento
que dio origen a este texto.  Los roles de autor se clasificaron
utilizando la taxonomía de roles de colaborador (CRediT;
https://credit.niso.org/) de la siguiente manera:  Edison
Achalma:   conceptualización, metodología, análisis
formal, investigación, recursos, curación de
datos, redacción, visualización, supervisión, administración del
proyecto}
\par{La correspondencia relativa a este artículo debe dirigirse a Edison
Achalma, Escuela Profesional de Economía, Universidad Nacional de San
Cristóbal de Huamanga, Portal Independencia N°
57, Ayacucho, AYA 5001, Perú, Email: \href{mailto:elmer.achalma.09@unsch.edu.pe}{elmer.achalma.09@unsch.edu.pe}}
}

\usepackage{pbalance}
% \usepackage{float}
\makeatletter
\let\oldtpt\ThreePartTable
\let\endoldtpt\endThreePartTable
\def\ThreePartTable{\@ifnextchar[\ThreePartTable@i \ThreePartTable@ii}
\def\ThreePartTable@i[#1]{\begin{figure}[!htbp]
\onecolumn
\begin{minipage}{0.485\textwidth}
\oldtpt[#1]
}
\def\ThreePartTable@ii{\begin{figure}[!htbp]
\onecolumn
\begin{minipage}{0.48\textwidth}
\oldtpt
}
\def\endThreePartTable{
\endoldtpt
\end{minipage}
\twocolumn
\end{figure}}
\makeatother


\makeatletter
\let\endoldlt\endlongtable		
\def\endlongtable{
\hline
\endoldlt}
\makeatother

\newenvironment{twocolumntable}% environment name
{% begin code
\begin{table*}[!htbp]%
\onecolumn%
}%
{%
\twocolumn%
\end{table*}%
}% end code

\urlstyle{same}



\usepackage{xurl}
\hypersetup{
  pdftitle={Institutos p\'{u}blicos de Ayacucho: sin luchas no hay victorias},
  pdfauthor={Autor}
}
\makeatletter
\@ifpackageloaded{tcolorbox}{}{\usepackage[skins,breakable]{tcolorbox}}
\@ifpackageloaded{fontawesome5}{}{\usepackage{fontawesome5}}
\definecolor{quarto-callout-color}{HTML}{909090}
\definecolor{quarto-callout-note-color}{HTML}{0758E5}
\definecolor{quarto-callout-important-color}{HTML}{CC1914}
\definecolor{quarto-callout-warning-color}{HTML}{EB9113}
\definecolor{quarto-callout-tip-color}{HTML}{00A047}
\definecolor{quarto-callout-caution-color}{HTML}{FC5300}
\definecolor{quarto-callout-color-frame}{HTML}{acacac}
\definecolor{quarto-callout-note-color-frame}{HTML}{4582ec}
\definecolor{quarto-callout-important-color-frame}{HTML}{d9534f}
\definecolor{quarto-callout-warning-color-frame}{HTML}{f0ad4e}
\definecolor{quarto-callout-tip-color-frame}{HTML}{02b875}
\definecolor{quarto-callout-caution-color-frame}{HTML}{fd7e14}
\makeatother
\makeatletter
\@ifpackageloaded{caption}{}{\usepackage{caption}}
\AtBeginDocument{%
\ifdefined\contentsname
  \renewcommand*\contentsname{Tabla de contenidos}
\else
  \newcommand\contentsname{Tabla de contenidos}
\fi
\ifdefined\listfigurename
  \renewcommand*\listfigurename{Lista de Figuras}
\else
  \newcommand\listfigurename{Lista de Figuras}
\fi
\ifdefined\listtablename
  \renewcommand*\listtablename{Lista de Tablas}
\else
  \newcommand\listtablename{Lista de Tablas}
\fi
\ifdefined\figurename
  \renewcommand*\figurename{Figura}
\else
  \newcommand\figurename{Figura}
\fi
\ifdefined\tablename
  \renewcommand*\tablename{Tabla}
\else
  \newcommand\tablename{Tabla}
\fi
}
\@ifpackageloaded{float}{}{\usepackage{float}}
\floatstyle{ruled}
\@ifundefined{c@chapter}{\newfloat{codelisting}{h}{lop}}{\newfloat{codelisting}{h}{lop}[chapter]}
\floatname{codelisting}{Listado}
\newcommand*\listoflistings{\listof{codelisting}{Lista de Listados}}
\makeatother
\makeatletter
\makeatother
\makeatletter
\@ifpackageloaded{caption}{}{\usepackage{caption}}
\@ifpackageloaded{subcaption}{}{\usepackage{subcaption}}
\makeatother

% From https://tex.stackexchange.com/a/645996/211326
%%% apa7 doesn't want to add appendix section titles in the toc
%%% let's make it do it
\makeatletter
\xpatchcmd{\appendix}
  {\par}
  {\addcontentsline{toc}{section}{\@currentlabelname}\par}
  {}{}
\makeatother

%% Disable longtable counter
%% https://tex.stackexchange.com/a/248395/211326

\usepackage{etoolbox}

\makeatletter
\patchcmd{\LT@caption}
  {\bgroup}
  {\bgroup\global\LTpatch@captiontrue}
  {}{}
\patchcmd{\longtable}
  {\par}
  {\par\global\LTpatch@captionfalse}
  {}{}
\apptocmd{\endlongtable}
  {\ifLTpatch@caption\else\addtocounter{table}{-1}\fi}
  {}{}
\newif\ifLTpatch@caption
\makeatother

\begin{document}

\maketitle


\hypertarget{toc}{}
\tableofcontents
\newpage
\section[Introduction]{Institutos públicos de Ayacucho: sin luchas no
hay victorias}


\setlength\LTleft{0pt}




\begin{tcolorbox}[enhanced jigsaw, arc=.35mm, bottomrule=.15mm, breakable, colback=white, colframe=quarto-callout-note-color-frame, left=2mm, leftrule=.75mm, opacityback=0, rightrule=.15mm, toprule=.15mm]
\begin{minipage}[t]{5.5mm}
\textcolor{quarto-callout-note-color}{\faInfo}
\end{minipage}%
\begin{minipage}[t]{\textwidth - 5.5mm}

\vspace{-3mm}\textbf{En pocas palabras}\vspace{3mm}

\begin{itemize}
\tightlist
\item
  \textbf{El problema.} En Ayacucho, 10 de los 16 institutos
  tecnológicos públicos del sector Educación no tienen licencia para
  funcionar con condiciones básicas de calidad. En cambio, todas las
  escuelas pedagógicas públicas de la región sí la tienen.
\item
  \textbf{Lo que sostengo.} La diferencia no se explica solo por falta
  de dinero. La ley exige a todos los institutos las mismas condiciones
  de calidad, pero solo garantiza recursos para infraestructura y
  equipamiento a las escuelas pedagógicas. Además, la reforma se aprobó
  sin presupuesto adicional, y el Gobierno Regional, que administra los
  institutos, les destina el 1,8 \% de su presupuesto.
\item
  \textbf{La propuesta.} Una coordinadora de estudiantes de institutos
  públicos de Ayacucho que convierta las demandas sueltas en pedidos con
  fundamento, destinatario, plazo e indicador.
\end{itemize}

\end{minipage}%
\end{tcolorbox}

Este texto parte de un pronunciamiento de estudiantes de institutos de
Ayacucho que circula con una consigna sencilla, \emph{sin luchas no hay
victorias}. El pronunciamiento denuncia institutos sin laboratorios, sin
equipos y sin locales adecuados, y llama a la unidad. Aquí intento darle
a esa denuncia el respaldo que merece: cifras oficiales, normas vigentes
y una explicación de por qué ocurre lo que ocurre.

En el camino separo tres tipos de afirmaciones. Cuando escribo «los
datos muestran», remito a una fuente que cualquiera puede consultar.
Cuando escribo «observé», hablo de lo que vi en visitas a algunos
institutos, sin pretender que valga para todos. Cuando escribo
«sostengo», es mi interpretación, y la someto a las objeciones que me
parecen más serias. El análisis se ocupa de los institutos y escuelas
públicos de educación superior de Ayacucho; la universidad y la
educación básica aparecen solo como contexto. Todas las cifras se
consultaron el 24 de octubre de 2026.

\section{La desigualdad empieza antes del
instituto}\label{la-desigualdad-empieza-antes-del-instituto}

Para muchos jóvenes de Ayacucho, el instituto público es la opción
realista de educación superior: no cobra pensión, dura tres años y suele
estar en la capital de la provincia. Esa función se entiende mejor
mirando quiénes llegan a la educación superior y quiénes no.

Los datos muestran que el acceso es estrecho y desigual. Entre 2014 y
2018, en promedio, solo tres de cada diez jóvenes de 15 a 19 años que
aprobaron el quinto de secundaria pasaron a la educación superior al año
siguiente. En el 20 \% de hogares con menor gasto por persona lo hizo el
21,8 \%; en el 20 \% con mayor gasto, el 52,7 \%. En Ayacucho la tasa
fue de 19,9 \%, una de las más bajas del país
(\citeproc{ref-ministeriodeeducacionPoliticaNacionalEducacion2020}{Ministerio
de Educación, 2020, pp. 26-27}). La cifra mide la \emph{transición
inmediata}, es decir, quién entra al año siguiente de terminar el
colegio; no cuenta a quienes entran más tarde.

La desigualdad también viene del colegio. En la prueba PISA de 2022, que
evalúa a estudiantes de quince años, solo el 33,8 \% de los peruanos
alcanzó al menos el nivel 2 en Matemática, el 47,4 \% en Ciencia y el
49,6 \% en Lectura
(\citeproc{ref-ministeriodeeducacionPeruPISA20222024}{Ministerio de
Educación, 2024, p. 46}). La Organización para la Cooperación y el
Desarrollo Económicos (OCDE) define el nivel 2 como la línea de base
desde la cual un estudiante empieza a mostrar las competencias
necesarias para participar en la vida como estudiante, trabajador y
ciudadano
(\citeproc{ref-ministeriodeeducacionPeruPISA20222024}{Ministerio de
Educación, 2024, p. 45}). Esa cifra describe a una muestra nacional de
estudiantes de quince años, no a los ingresantes de un instituto
concreto. Sirve para algo más modesto: recordar que los institutos
reciben a jóvenes de un sistema escolar con déficits conocidos y que
nivelarlos exige recursos.

¿Qué ofrece el Estado a esos jóvenes? Para el semestre 2025-I, 343
institutos tecnológicos públicos en proceso de admisión ofrecieron 48
196 vacantes gratuitas en todo el país
(\citeproc{ref-ministeriodeeducacionInstitutosPublicosOfrecen2025}{Ministerio
de Educación, 2025b}). Aun así, el sector privado concentra la
matrícula: según el Censo Educativo 2025, de cada cuatro estudiantes de
educación superior tecnológica, casi tres están en un instituto privado
(385 626 frente a 138 211), y en Ayacucho la proporción es la misma, 9
390 frente a 3 341
(\citeproc{ref-ministeriodeeducacionPadronInstitucionesEducativas2025}{Ministerio
de Educación, 2025a}).

La región no es una excepción. La expansión de la educación superior
latinoamericana en lo que va del siglo se apoyó en buena parte en
instituciones privadas, llevó a ellas a estudiantes de menores ingresos
y produjo una calidad muy desigual entre instituciones
(\citeproc{ref-ferreyraCrossroadsHigherEducation2017}{Ferreyra et~al.,
2017}). Para los programas técnicos cortos, como los de un instituto, la
evidencia regional indica que suelen mejorar los ingresos de sus
egresados frente a quienes solo terminaron la secundaria, aunque esos
resultados varían mucho de un programa a otro
(\citeproc{ref-ferreyraFastTrackNew2021}{Ferreyra et~al., 2021}). En el
Perú, además, el subempleo de los egresados de educación superior es un
problema persistente, como documentan Yamada y Lavado
(\citeproc{ref-yamadaEducacionSuperiorEmpleo2017}{2017}). La calidad del
instituto, entonces, no es un detalle: decide si el título abre una
puerta o no.

\section{Los institutos públicos de Ayacucho, uno por
uno}\label{los-institutos-puxfablicos-de-ayacucho-uno-por-uno}

Ayacucho tiene 16 institutos tecnológicos públicos del sector Educación,
uno o más en cada una de sus once provincias, además de la escuela
técnica de la Policía Nacional en Huamanga. Tiene también cinco escuelas
pedagógicas públicas y dos escuelas de formación artística públicas
(\citeproc{ref-ministeriodeeducacionPadronInstitucionesEducativas2025}{Ministerio
de Educación, 2025a}). La Tabla~\ref{tbl-ayacucho} muestra los
tecnológicos, con su matrícula y si tienen licencia.

\begin{twocolumntable}

\begin{longtable}[]{@{}llrl@{}}
\caption{Institutos tecnológicos públicos del sector Educación en
Ayacucho}\label{tbl-ayacucho}\tabularnewline
\toprule\noalign{}
Instituto & Provincia & Matrícula & Licencia \\
\midrule\noalign{}
\endfirsthead
\toprule\noalign{}
Instituto & Provincia & Matrícula & Licencia \\
\midrule\noalign{}
\endhead
\bottomrule\noalign{}
\endlastfoot
Víctor Álvarez Huapaya & Huamanga & 545 & Sí \\
Manuel Antonio Hierro Pozo & Huamanga & 430 & Sí \\
Huanta & Huanta & 384 & Sí \\
José María Arguedas & Lucanas (Puquio) & 358 & No \\
César Augusto Guardia Mayorga & Parinacochas (Coracora) & 352 & No \\
Perú Corea del Sur & Víctor Fajardo (Huancapi) & 146 & Sí \\
San Miguel & La Mar & 145 & No \\
Federico Gonzales Cabezudo & Huanca Sancos & 137 & No \\
Virgen del Rosario & Lucanas (Laramate) & 119 & No \\
Aucará & Lucanas & 110 & Sí \\
Chipao & Lucanas & 89 & No \\
Paucar del Sara Sara & Paucar del Sara Sara (Pausa) & 79 & Sí \\
Santo Domingo de Guzmán & Sucre (Querobamba) & 71 & No \\
San Juan & Vilcas Huamán & 71 & No \\
Los Morochucos & Cangallo & 62 & No \\
Hatun Soras & Sucre (Soras) & 51 & No \\
\textbf{Total} & & \textbf{3 149} & \textbf{6 de 16} \\
\end{longtable}

\noindent \emph{Nota.} Matrícula:~padrón del Censo Educativo 2025 del
Ministerio de Educación. Licencia: lista oficial de institutos
licenciados del Ministerio de Educación, con corte al 23 de junio de
2026. No incluye la escuela técnica de la Policía Nacional, que depende
de otro sector.

\end{twocolumntable}

Los datos muestran tres rasgos. El primero es la escala: diez de los
dieciséis tienen menos de 150 estudiantes, y entre todos suman 349
docentes, unos nueve estudiantes por docente. El segundo es la brecha de
licencia: 6 de 16 tecnológicos públicos tienen licencia, mientras que
las cinco escuelas pedagógicas públicas ya la obtuvieron. Ayacucho fue
la primera región del país en licenciar a todas sus escuelas pedagógicas
públicas, con la de Puquio en mayo de 2026
(\citeproc{ref-ministeriodeeducacionRegionAyacuchoLogra2026}{Ministerio
de Educación, 2026b}). El tercero es que el tamaño no lo explica todo:
Paucar del Sara Sara, con 79 estudiantes, tiene licencia, y José María
Arguedas, con 358, no.

En el país, el panorama es parecido. Hasta junio de 2026 el Ministerio
de Educación había licenciado 246 institutos desde 2018, y solo 51 eran
públicos
(\citeproc{ref-ministeriodeeducacionListaInstitutosLicenciados2026}{Ministerio
de Educación, 2026a}). El licenciamiento es el procedimiento que
verifica que un instituto cumple condiciones básicas de calidad para
funcionar, entre ellas una infraestructura y un equipamiento adecuados
(\citeproc{ref-ministeriodeeducacionLeyNdeg305122016}{Ley N.° 30512,
2016}, arts. 24-25).

\subsection{Lo que observé, y sus
límites}\label{lo-que-observuxe9-y-sus-luxedmites}

En visitas recientes a institutos del sur de Ayacucho observé locales de
adobe con techos que dejan pasar el agua, una carrera de Enfermería
Técnica con un solo microscopio para todos sus estudiantes y domos de
emergencia guardados en un almacén, a la espera de que alguien los arme.
Según lo que pude conocer en esas visitas, varios ya tienen expediente
técnico para su nuevo local y esperan financiamiento. Son observaciones
de unos pocos institutos, no una muestra: no permiten afirmar que todos
estén igual, y el estado de cada proyecto debe comprobarse en el banco
de inversiones del Estado. Ese es justamente uno de los primeros
trabajos que propongo más adelante.

\section{Presupuesto: cuánto llega y a
dónde}\label{presupuesto-cuuxe1nto-llega-y-a-duxf3nde}

Para hablar de dinero con precisión hay que distinguir tres medidas. El
\emph{gasto por estudiante} compara cuánto invierte un país en cada
estudiante, en relación con su ingreso por habitante. El
\emph{presupuesto institucional modificado} (PIM) es lo que se asigna en
el año, después de los cambios. El \emph{devengado} es lo que
efectivamente se gastó. Y dentro del gasto conviene separar el de
funcionamiento (sueldos, servicios, materiales) del de inversión (obras
y equipos nuevos).

\textbf{Poco en comparación con los vecinos.} En 2017, último año con
dato para el Perú, el gasto público por estudiante de educación
terciaria fue el 10,9 \% del producto bruto interno (PBI) por habitante,
frente a 20,3 \% en Chile y 20,4 \% en Colombia
(\citeproc{ref-bancomundialGastoPublicoEstudiante2026}{Banco Mundial,
2026}). Por cada estudiante de educación superior, el Estado peruano
invertía, en proporción a su economía, algo más de la mitad que sus
vecinos.

\textbf{Mejor repartido que antes, pero todavía desigual.} La
Tabla~\ref{tbl-gasto} muestra cómo se distribuye el presupuesto público
de la educación superior y técnico-productiva.

\begin{twocolumntable}

\begin{longtable}[]{@{}lrrr@{}}
\caption{Distribución del presupuesto público de la educación superior y
técnico-productiva}\label{tbl-gasto}\tabularnewline
\toprule\noalign{}
Año y medida & Universidad & Institutos & CETPRO \\
\midrule\noalign{}
\endfirsthead
\toprule\noalign{}
Año y medida & Universidad & Institutos & CETPRO \\
\midrule\noalign{}
\endhead
\bottomrule\noalign{}
\endlastfoot
2019, asignado (PIM) & 73,4 \% & 20,4 \% & 6,2 \% \\
2025, asignado (PIM) & 69,0 \% & 25,5 \% & 5,5 \% \\
2025, gastado (devengado) & 68,2 \% & 25,9 \% & 5,9 \% \\
\end{longtable}

\noindent \emph{Nota.} Elaboración~propia con la Consulta Amigable del
MEF (Ministerio de Economía y Finanzas, 2026), todos los niveles de
gobierno; el dato de 2019 coincide con el de la política nacional
(Ministerio de Educación, 2020, pp.~46-47). Institutos: institutos y
escuelas, grupo funcional 0108, que reúne tecnológicos, pedagógicos y
artísticos. CETPRO: centros de educación técnico-productiva. No incluye
el grupo de infraestructura y equipamiento, que el sistema no reparte
por tipo de institución. Montos nominales.

\end{twocolumntable}

Los datos muestran una mejora que conviene reconocer. Lo asignado a
institutos y escuelas pasó de S/ 1 119 millones en 2019 a S/ 2 045
millones en 2025, y su participación subió de una quinta a una cuarta
parte
(\citeproc{ref-ministeriodeeconomiayfinanzasConsultaAmigable2026}{Ministerio
de Economía y Finanzas, 2026}). La universidad, con todo, sigue
recibiendo casi tres veces más. Hay un límite importante: el sistema de
presupuesto no permite separar lo que va a tecnológicos, a pedagógicos y
a escuelas artísticas, ni seguir por región el gasto en infraestructura.
La propia política nacional lo reconoce
(\citeproc{ref-ministeriodeeducacionPoliticaNacionalEducacion2020}{Ministerio
de Educación, 2020, p. 120}).

\textbf{En Ayacucho.} En 2025, el Gobierno Regional de Ayacucho tuvo S/
49,9 millones para educación superior no universitaria y gastó S/ 49,5
millones, el 99 \%. Ese monto equivale al 1,8 \% de su presupuesto
total, de S/ 2 788 millones, y es más del doble de los S/ 24 millones de
2019
(\citeproc{ref-ministeriodeeconomiayfinanzasConsultaAmigable2026}{Ministerio
de Economía y Finanzas, 2026}). Casi todo lo que el Estado gasta en
institutos dentro de la región pasa por el Gobierno Regional. Si se
divide ese monto entre los cerca de 6 000 estudiantes de los
tecnológicos, pedagógicos y artísticos públicos, salen unos S/ 8 350 por
estudiante al año. Es un cálculo aproximado: incluye las remuneraciones
de docentes y personal, y no solo laboratorios u obras.

\section{Lo que dice la ley y quién
responde}\label{lo-que-dice-la-ley-y-quiuxe9n-responde}

\subsection{Una cronología breve}\label{una-cronologuxeda-breve}

La Tabla~\ref{tbl-normas} resume las normas que importan para este
análisis y su situación a octubre de 2026.

\begin{twocolumntable}

\begin{longtable}[]{@{}
  >{\raggedright\arraybackslash}p{(\linewidth - 2\tabcolsep) * \real{0.5000}}
  >{\raggedright\arraybackslash}p{(\linewidth - 2\tabcolsep) * \real{0.5000}}@{}}
\caption{Normas sobre los institutos públicos, de 1996 a
2022}\label{tbl-normas}\tabularnewline
\toprule\noalign{}
\begin{minipage}[b]{\linewidth}\raggedright
Norma
\end{minipage} & \begin{minipage}[b]{\linewidth}\raggedright
Qué importa hoy
\end{minipage} \\
\midrule\noalign{}
\endfirsthead
\toprule\noalign{}
\begin{minipage}[b]{\linewidth}\raggedright
Norma
\end{minipage} & \begin{minipage}[b]{\linewidth}\raggedright
Qué importa hoy
\end{minipage} \\
\midrule\noalign{}
\endhead
\bottomrule\noalign{}
\endlastfoot
Decreto Legislativo N.° 882
(\citeproc{ref-peruDecretoLegislativo8821996}{1996}) & Permitió la
educación privada con fines de lucro; derogado en 2016 para institutos y
escuelas \\
Ley N.° 27867 (\citeproc{ref-peruLeyOrganicaGobiernos2002}{2002}),
Orgánica de Gobiernos Regionales & El gobierno regional gestiona la
educación superior no universitaria y su infraestructura \\
Ley N.° 30512
(\citeproc{ref-ministeriodeeducacionLeyNdeg305122016}{2016}), de
Institutos y Escuelas & Licenciamiento, bienestar estudiantil, 10 \% del
canon; reforma sin recursos adicionales \\
Decreto Supremo N.° 010-2017-MINEDU
(\citeproc{ref-peruReglamentoLey305122017}{2017}) & Reglamento de la Ley
30512 \\
Decreto de Urgencia N.° 017-2020
(\citeproc{ref-peruDecretoUrgencia0172020}{2020}) & Modifica el
licenciamiento y permite formas societarias a los institutos privados \\
Ley N.° 31653 (\citeproc{ref-peruLey316532022}{2022}) & Permite que las
escuelas pedagógicas y tecnológicas públicas sean unidades ejecutoras \\
\end{longtable}

\noindent \emph{Nota.} Elaboración~propia a partir de los textos
publicados en el diario oficial El Peruano.

\end{twocolumntable}

El Decreto Legislativo 882 reconoció el derecho a conducir instituciones
educativas «con o sin finalidad lucrativa» y les extendió las garantías
de la inversión privada
(\citeproc{ref-peruDecretoLegislativo8821996}{Decreto Legislativo N.°
882, 1996}, arts. 2 y 7). La Ley 30512 lo derogó en 2016 para institutos
y escuelas (\citeproc{ref-ministeriodeeducacionLeyNdeg305122016}{Ley N.°
30512, 2016}, décima tercera disposición complementaria final), pero la
gestión privada con fines de lucro sigue permitida bajo las formas del
derecho común o societario
(\citeproc{ref-peruDecretoUrgencia0172020}{Decreto de Urgencia N.°
017-2020, 2020}, que modifica el art. 9 de la Ley 30512). Sostengo que
ese decreto abrió la puerta al crecimiento privado, pero no basta para
explicar la precariedad de lo público: que lo privado creciera no
obligaba a que lo público se quedara atrás. Esa segunda parte se explica
por cómo se financian los institutos públicos, que es lo que sigue.

\subsection{Obligación, insuficiencia y
vacío}\label{obligaciuxf3n-insuficiencia-y-vacuxedo}

La Ley 30512 trata de manera distinta a dos tipos de instituciones
públicas, y conviene separar tres situaciones.

\textbf{Donde hay una obligación expresa.} La ley obliga al Ministerio
de Educación y a los gobiernos regionales a considerar en sus
presupuestos, «bajo responsabilidad», los recursos para la
infraestructura, el mobiliario y el equipamiento de las escuelas de
educación superior pedagógica públicas
(\citeproc{ref-ministeriodeeducacionLeyNdeg305122016}{Ley N.° 30512,
2016}, arts. 61 y 63). También manda que cada gobierno regional entregue
el 10 \% del canon que recibe a los institutos y escuelas de su
circunscripción, solo para invertir en investigación científica y
tecnológica y en su infraestructura
(\citeproc{ref-ministeriodeeducacionLeyNdeg305122016}{Ley N.° 30512,
2016}, tercera disposición complementaria modificatoria). Si estas
obligaciones no se cumplen, hay un incumplimiento que se puede exigir.

\textbf{Donde la obligación es general y el dinero no está asegurado.}
Para los institutos tecnológicos públicos, la ley dice que el gobierno
regional provee el servicio y los gestiona
(\citeproc{ref-ministeriodeeducacionLeyNdeg305122016}{Ley N.° 30512,
2016}, art. 57), y la Ley Orgánica de Gobiernos Regionales le encarga
diseñar e implementar las políticas de infraestructura y equipamiento
(\citeproc{ref-peruLeyOrganicaGobiernos2002}{Ley N.° 27867, 2002}, art.
47, inc. m). Pero no existe una cláusula equivalente a la de las
escuelas pedagógicas. El reglamento, incluso, sujeta la investigación de
los institutos públicos «a la disponibilidad presupuestal»
(\citeproc{ref-peruReglamentoLey305122017}{Decreto Supremo N.°
010-2017-MINEDU, 2017}, art. 49.2). Aquí el problema no es un
incumplimiento, sino una garantía insuficiente: hay a quién exigir, pero
no hay un piso de recursos.

\textbf{Donde no hay norma.} La ley exige a cada instituto público una
unidad de bienestar y empleabilidad, con orientación profesional,
tutoría y consejería, y un comité de defensa del estudiante para casos
de acoso o discriminación
(\citeproc{ref-ministeriodeeducacionLeyNdeg305122016}{Ley N.° 30512,
2016}, art. 29, inc. f). No menciona un área de psicología. Pedirla es
legítimo, pero es una propuesta de política, no el cumplimiento de una
obligación.

A esto se suma la cláusula de financiamiento. La misma ley dispuso que,
después de 2016, su aplicación se financiara con el presupuesto de los
institutos y del Ministerio de Educación, «sin demandar recursos
adicionales al Tesoro Público»
(\citeproc{ref-ministeriodeeducacionLeyNdeg305122016}{Ley N.° 30512,
2016}, décima cuarta disposición complementaria transitoria). Se
exigieron nuevas condiciones de calidad y se dispuso que se pagaran con
lo que ya había.

\subsection{Quién decide, quién paga, quién
ejecuta}\label{quiuxe9n-decide-quiuxe9n-paga-quiuxe9n-ejecuta}

\begin{itemize}
\tightlist
\item
  \textbf{El Gobierno Regional de Ayacucho}, por medio de su Dirección
  Regional de Educación, gestiona los institutos públicos, administra su
  presupuesto y decide qué obras y equipos se financian. Mientras no
  funcione el organismo nacional previsto por la ley para gestionar los
  tecnológicos (Educatec), los gobiernos regionales siguen a cargo
  (\citeproc{ref-ministeriodeeducacionLeyNdeg305122016}{Ley N.° 30512,
  2016}, séptima disposición complementaria transitoria).
\item
  \textbf{El Ministerio de Educación} dicta la política, otorga o niega
  la licencia y puede transferir recursos.
\item
  \textbf{El Congreso} puede modificar la ley, por ejemplo para extender
  la garantía del artículo 61 a los tecnológicos, y sus representantes
  pueden fiscalizar y pedir información. No ejecutan el presupuesto
  regional.
\item
  \textbf{Las municipalidades} no tienen a su cargo los institutos, pero
  pueden apoyar con transporte o locales mediante convenios.
\end{itemize}

\section{Por qué ocurre: una lectura de economía
política}\label{por-quuxe9-ocurre-una-lectura-de-economuxeda-poluxedtica}

José Carlos Mariátegui escribió que «el problema de la enseñanza no
puede ser bien comprendido en nuestro tiempo, si no es considerado un
problema económico y como un problema social»
(\citeproc{ref-mariateguiSieteEnsayosInterpretacion2007}{Mariátegui,
1928/2007, p. 129}). Leído así, la pregunta no es solo cuánto dinero
falta, sino por qué el sistema produce esa falta. Propongo el siguiente
mecanismo, que es una interpretación apoyada en los datos anteriores:

\begin{enumerate}
\def\labelenumi{\arabic{enumi}.}
\tightlist
\item
  La ley fija condiciones de calidad iguales para todos los institutos y
  condiciona su funcionamiento a la licencia.
\item
  El costo de cumplirlas recae en el gobierno regional, que gestiona los
  institutos.
\item
  Para los tecnológicos no hay un piso legal de recursos, y la reforma
  se aprobó sin presupuesto adicional.
\item
  El gobierno regional reparte un presupuesto limitado entre salud,
  carreteras, colegios y muchas otras demandas. Los institutos son
  pequeños, están dispersos en provincias y tienen poca voz organizada:
  reciben el 1,8 \%.
\item
  El instituto que no reúne las condiciones no obtiene la licencia, y el
  joven que no puede pagar un privado se queda con menos opciones.
\end{enumerate}

La comparación con las escuelas pedagógicas apoya este mecanismo. Están
en la misma región, dependen del mismo gobierno regional y tienen la
garantía del artículo 61: todas están licenciadas. No es una prueba de
causalidad, porque las escuelas pedagógicas también difieren en tamaño y
en historia, pero es un contraste difícil de ignorar.

Sostengo, además, que aprobar normas sin presupuesto responde a un
incentivo político. Quien aprueba una reforma gana el reconocimiento de
inmediato, mientras que el costo de cumplirla se traslada a otras
entidades y a otros años. Pierre Bourdieu describió el neoliberalismo
como una política que, en nombre de la eficiencia del mercado, restringe
el gasto público y funciona como «un programa de destrucción metódica de
los colectivos»
(\citeproc{ref-bourdieuEsenciaNeoliberalismo101}{Bourdieu, 1998, p. 3},
traducción propia). No hace falta suponer una intención deliberada para
ver el efecto: un sistema que deja que cada instituto resuelva solo sus
carencias debilita justamente lo que podría corregirlas, la acción
conjunta de quienes estudian y enseñan en ellos.

\subsection{Otras explicaciones que conviene
considerar}\label{otras-explicaciones-que-conviene-considerar}

\textbf{La capacidad de ejecución.} Se dice a menudo que el problema no
es la falta de dinero, sino que no se gasta. En el grupo de educación
superior no universitaria, el Gobierno Regional ejecutó el 99 \% de lo
asignado en 2025, de modo que ahí la ejecución no es el cuello de
botella. Queda abierta la ejecución de las obras, que se registran
aparte y que esta revisión no pudo seguir.

\textbf{La dispersión territorial.} Dieciséis institutos para 3 149
estudiantes es una red cara de mantener, y alguien podría proponer
concentrarlos. Pero esa cercanía es justamente lo que permite estudiar a
quien no puede mudarse a Huamanga, en una región donde solo uno de cada
cinco egresados de secundaria pasa de inmediato a la educación superior.
La discusión de fondo es cómo financiar una red dispersa, no si debe
existir.

\textbf{La gestión de cada instituto.} La diferencia entre Paucar del
Sara Sara, licenciado con 79 estudiantes, y otros más grandes sin
licencia sugiere que la gestión importa. Eso no contradice el argumento:
la gestión puede aprovechar mejor los recursos, pero no reemplaza la
falta de un piso de financiamiento.

\section{Organizarse para exigir}\label{organizarse-para-exigir}

Si el problema es que cada instituto enfrenta solo una carencia que es
común, la respuesta lógica es actuar en conjunto. El pronunciamiento
propone una \textbf{Coordinadora de estudiantes de institutos de
Ayacucho}. Todavía es una propuesta, no una organización constituida.
Para que funcione, sugiero empezar así:

\begin{itemize}
\tightlist
\item
  \textbf{Representación plural.} Un delegado por cada consejo
  estudiantil de los 16 tecnológicos, las 5 escuelas pedagógicas y las 2
  escuelas artísticas públicas, para que estén las once provincias. Las
  decisiones, por mayoría de delegados y con actas públicas.
\item
  \textbf{Independencia.} Es una iniciativa ciudadana de estudiantes,
  sin respaldo partidario ni institucional.
\item
  \textbf{Primer paso, en 30 días.} Un diagnóstico común con la misma
  ficha en cada instituto (equipos, locales, conectividad, bienestar), y
  solicitudes de acceso a la información pública al Gobierno Regional
  sobre los proyectos de inversión de cada instituto y sobre el 10 \%
  del canon. La entidad debe responder en un plazo de diez días hábiles
  (\citeproc{ref-peruTUOLeyTransparencia2019}{Decreto Supremo N.°
  021-2019-JUS, 2019}, art. 11, inc. b).
\item
  \textbf{Segundo paso, en 60 días.} Con el diagnóstico y las
  respuestas, un pliego regional único y el pedido de una audiencia
  pública con el Gobierno Regional y la Dirección Regional de Educación.
\item
  \textbf{Seguimiento.} Un registro público de cada demanda, con lo
  pedido, la respuesta y el avance.
\end{itemize}

Para ordenar las prioridades, propongo cuatro criterios: la seguridad de
los estudiantes (un techo que se cae va primero), el número de
estudiantes afectados, el efecto sobre la licencia del instituto y el
costo.

\section{Agenda regional de demandas}\label{agenda-regional-de-demandas}

La Tabla~\ref{tbl-demandas} convierte el pliego del pronunciamiento en
una matriz de seguimiento. Distingue las demandas inmediatas, que se
resuelven con presupuesto, de las estructurales, que requieren cambiar
una norma. Todas están hoy en estado \emph{pendiente}: documentarlas es
la primera tarea de la coordinadora.

\begin{twocolumntable}

\begin{longtable}[]{@{}
  >{\raggedright\arraybackslash}p{(\linewidth - 4\tabcolsep) * \real{0.4000}}
  >{\raggedright\arraybackslash}p{(\linewidth - 4\tabcolsep) * \real{0.3250}}
  >{\raggedright\arraybackslash}p{(\linewidth - 4\tabcolsep) * \real{0.2750}}@{}}
\caption{Matriz regional de seguimiento de las demandas
estudiantiles}\label{tbl-demandas}\tabularnewline
\toprule\noalign{}
\begin{minipage}[b]{\linewidth}\raggedright
Demanda y acción solicitada
\end{minipage} & \begin{minipage}[b]{\linewidth}\raggedright
Fundamento
\end{minipage} & \begin{minipage}[b]{\linewidth}\raggedright
Responsable e indicador
\end{minipage} \\
\midrule\noalign{}
\endfirsthead
\toprule\noalign{}
\begin{minipage}[b]{\linewidth}\raggedright
Demanda y acción solicitada
\end{minipage} & \begin{minipage}[b]{\linewidth}\raggedright
Fundamento
\end{minipage} & \begin{minipage}[b]{\linewidth}\raggedright
Responsable e indicador
\end{minipage} \\
\midrule\noalign{}
\endhead
\bottomrule\noalign{}
\endlastfoot
Laboratorios, equipos y maquinaria de las carreras, como Enfermería
Técnica y las agropecuarias (inmediata). Inventario por instituto y
pedido presupuestal para 2027 & Condiciones básicas de calidad (Ley
30512, arts. 24-25); gestión regional (art. 57; Ley 27867, art. 47 m) &
GR y DREA, con apoyo del Minedu. Equipos operativos por carrera \\
Locales nuevos o rehabilitados (inmediata). Lista de proyectos con
código, costo y avance & Ley 27867, art. 47 m; Ley 30512, art. 57 & GR.
Avance físico y financiero de cada obra \\
Licencia para los 10 tecnológicos públicos que no la tienen (inmediata).
Plan con brechas y costo por instituto & Ley 30512, arts. 24-25 & GR y
Minedu. Institutos licenciados \\
Bibliotecas, salas de cómputo e internet en todo el local (inmediata) &
Condiciones básicas de calidad (Ley 30512, arts. 24-25) & GR. Conexión
operativa en cada instituto \\
Unidad de bienestar con consejería y comité de defensa; atención de
salud mental (inmediata) & Ley 30512, art. 29 f; el área de psicología
es una propuesta, no una obligación & Dirección del instituto y GR.
Unidad funcionando con personal \\
Comedor, residencia y movilidad (inmediata) & Propuesta de política;
acceso equitativo como objetivo de la política nacional & GR y
municipalidades. Convenios firmados y estudiantes atendidos \\
Personal de vigilancia (inmediata) & Gasto de funcionamiento del
presupuesto regional & GR. Plazas cubiertas \\
El 10 \% del canon para investigación en institutos (estructural) & Ley
30512, tercera disposición complementaria modificatoria & GR. Montos
transferidos por año \\
Garantía de recursos del art. 61 también para los tecnológicos públicos
(estructural) & Reforma legal: hoy solo rige para las escuelas
pedagógicas & Congreso. Proyecto de ley presentado y debatido \\
\end{longtable}

\noindent \emph{Nota.} Elaboración~propia a partir del pronunciamiento
estudiantil, la Ley N.° 30512 (2016) y la Ley N.° 27867 (2002). GR:
Gobierno Regional de Ayacucho; DREA: Dirección Regional de Educación de
Ayacucho.

\end{twocolumntable}

\section{Conclusiones}\label{conclusiones}

Los datos permiten afirmar tres cosas. El Estado peruano invierte por
estudiante de educación superior algo más de la mitad que Chile o
Colombia. Los institutos reciben hoy una cuarta parte del presupuesto
público de la educación superior, más que en 2019, pero mucho menos que
la universidad. Y en Ayacucho, 10 de los 16 tecnológicos públicos no
tienen licencia, mientras que todas las escuelas pedagógicas públicas
sí.

La explicación que propongo, y que dejo abierta a discusión, es que la
ley exige a los tecnológicos la misma calidad que a las escuelas
pedagógicas sin garantizarles los mismos recursos, y que la reforma se
aprobó sin presupuesto adicional. Esa brecha entre lo que se exige y lo
que se financia la paga quien no puede elegir un instituto privado.

La coordinadora es todavía una propuesta. Su primer paso es concreto:
reunir a los delegados, levantar un diagnóstico común y pedir
información pública. Con ese diagnóstico en la mano, las demandas dejan
de ser un reclamo disperso y se vuelven un pliego con fundamento,
responsable y plazo.

\textbf{¡Por la unidad de los estudiantes de institutos de la región de
Ayacucho! ¡Sin luchas no hay victorias!}

\section{Referencias}\label{referencias}

\protect\phantomsection\label{refs}
\begin{CSLReferences}{1}{0}
\bibitem[\citeproctext]{ref-bancomundialGastoPublicoEstudiante2026}
Banco Mundial. (2026). \emph{Gasto público por estudiante, nivel
terciario (\% del {PIB} per cápita), indicador {SE.XPD.TERT.PC.ZS}}
{[}Conjunto de datos{]}. Instituto de Estadística de la UNESCO.
\url{https://datos.bancomundial.org/indicador/SE.XPD.TERT.PC.ZS?locations=PE-CL-CO}

\bibitem[\citeproctext]{ref-bourdieuEsenciaNeoliberalismo101}
Bourdieu, P. (1998). L'essence du néolibéralisme. \emph{Le Monde
diplomatique}, 3.
\url{https://www.monde-diplomatique.fr/1998/03/BOURDIEU/3609}

\bibitem[\citeproctext]{ref-peruDecretoUrgencia0172020}
Decreto de Urgencia N.° 017-2020. (2020). \emph{Decreto de urgencia que
establece medidas para el fortalecimiento de la gestión y el
licenciamiento de los institutos y escuelas de educación superior, en el
marco de la Ley N.° 30512}. Diario Oficial El Peruano.
\url{https://static.legis.pe/wp-content/uploads/2020/01/Decreto-Urgencia-17-2020-LP.pdf}

\bibitem[\citeproctext]{ref-peruDecretoLegislativo8821996}
Decreto Legislativo N.° 882. (1996). \emph{Ley de Promoción de la
Inversión en la Educación}. Diario Oficial El Peruano.
\url{https://www.gob.pe/institucion/minedu/normas-legales/118067-882}

\bibitem[\citeproctext]{ref-peruReglamentoLey305122017}
Decreto Supremo N.° 010-2017-MINEDU. (2017). \emph{Reglamento de la Ley
N.° 30512}. Diario Oficial El Peruano.
\url{https://www.minedu.gob.pe/superiorpedagogica/download/99/}

\bibitem[\citeproctext]{ref-peruTUOLeyTransparencia2019}
Decreto Supremo N.° 021-2019-JUS. (2019). \emph{Texto Único Ordenado de
la Ley N.° 27806, Ley de Transparencia y Acceso a la Información
Pública}. Diario Oficial El Peruano.

\bibitem[\citeproctext]{ref-ferreyraCrossroadsHigherEducation2017}
Ferreyra, M. M., Avitabile, C., Botero Álvarez, J., Haimovich Paz, F.,
\& Urzúa, S. (2017). \emph{At a crossroads: Higher education in {Latin
America} and the {Caribbean}}. Banco Mundial.
\url{https://doi.org/10.1596/978-1-4648-1014-5}

\bibitem[\citeproctext]{ref-ferreyraFastTrackNew2021}
Ferreyra, M. M., Dinarte Díaz, L., Urzúa, S., \& Bassi, M. (2021).
\emph{The fast track to new skills: Short-cycle higher education
programs in {Latin America} and the {Caribbean}}. Banco Mundial.
\url{https://doi.org/10.1596/978-1-4648-1706-9}

\bibitem[\citeproctext]{ref-peruLeyOrganicaGobiernos2002}
Ley N.° 27867. (2002). \emph{Ley Orgánica de Gobiernos Regionales}.
Diario Oficial El Peruano.
\url{https://cdn.www.gob.pe/uploads/document/file/9590889/7851250-ley-n-27867-ley-organica-de-gobiernos-regionales.pdf}

\bibitem[\citeproctext]{ref-ministeriodeeducacionLeyNdeg305122016}
Ley N.° 30512. (2016). \emph{Ley de Institutos y Escuelas de Educación
Superior y de la Carrera Pública de sus Docentes}. Diario Oficial El
Peruano.
\url{https://www.gob.pe/institucion/minedu/normas-legales/118500-30512}

\bibitem[\citeproctext]{ref-peruLey316532022}
Ley N.° 31653. (2022). \emph{Ley que modifica la Ley 30512, Ley de
Institutos y Escuelas de Educación Superior y de la Carrera Pública de
sus Docentes}. Diario Oficial El Peruano.

\bibitem[\citeproctext]{ref-mariateguiSieteEnsayosInterpretacion2007}
Mariátegui, J. C. (2007). \emph{Siete ensayos de interpretación de la
realidad peruana} (3.ª ed.). Fundación Biblioteca Ayacucho. (Obra
original publicada en 1928)

\bibitem[\citeproctext]{ref-ministeriodeeconomiayfinanzasConsultaAmigable2026}
Ministerio de Economía y Finanzas. (2026). \emph{Consulta amigable:
consulta de ejecución del gasto} {[}Base de datos{]}. Portal de
Transparencia Económica.
\url{https://apps5.mineco.gob.pe/transparencia/Navegador/default.aspx}

\bibitem[\citeproctext]{ref-ministeriodeeducacionPoliticaNacionalEducacion2020}
Ministerio de Educación. (2020). \emph{Política Nacional de Educación
Superior y Técnico-Productiva}. Ministerio de Educación.
\url{https://siteal.iiep.unesco.org/sites/default/files/sit_accion_files/politica_nacional_de_educacion_superior_y_tecnico-productiva.pdf}

\bibitem[\citeproctext]{ref-ministeriodeeducacionPeruPISA20222024}
Ministerio de Educación. (2024). \emph{El {Perú} en {PISA} 2022: informe
nacional de resultados}. Oficina de Medición de la Calidad de los
Aprendizajes. \url{https://hdl.handle.net/20.500.12799/10199}

\bibitem[\citeproctext]{ref-ministeriodeeducacionPadronInstitucionesEducativas2025}
Ministerio de Educación. (2025a). \emph{Padrón de instituciones
educativas, {Censo Educativo} 2025} {[}Conjunto de datos{]}. Estadística
de la Calidad Educativa (ESCALE).
\url{https://escale.minedu.gob.pe/documents/inicio/Bases\%20de\%20datos/A.\%20Censo\%20Educativo/2025/00_Padron.zip}

\bibitem[\citeproctext]{ref-ministeriodeeducacionInstitutosPublicosOfrecen2025}
Ministerio de Educación. (2025b, abril 6). \emph{Institutos públicos
ofrecen 48 000 vacantes para estudiar gratis una carrera técnica}.
\url{https://www.gob.pe/institucion/minedu/noticias/1140949}

\bibitem[\citeproctext]{ref-ministeriodeeducacionListaInstitutosLicenciados2026}
Ministerio de Educación. (2026a). \emph{Lista de 246 institutos de
educación superior licenciados por el {MINEDU}}. Ministerio de
Educación.
\url{https://www.gob.pe/institucion/minedu/informes-publicaciones/4718331}

\bibitem[\citeproctext]{ref-ministeriodeeducacionRegionAyacuchoLogra2026}
Ministerio de Educación. (2026b, mayo 14). \emph{Región Ayacucho logra
el licenciamiento del 100 \% de sus escuelas pedagógicas públicas (Nota
de prensa N.° 026-2026)}.
\url{https://www.gob.pe/institucion/minedu/noticias/1390506}

\bibitem[\citeproctext]{ref-yamadaEducacionSuperiorEmpleo2017}
Yamada, G., \& Lavado, P. (Eds.). (2017). \emph{Educación superior y
empleo en el {Perú}: una brecha persistente}. Universidad del Pacífico.

\end{CSLReferences}






\end{document}
